%(BEGIN_QUESTION)
% Copyright 2015, Tony R. Kuphaldt, released under the Creative Commons Attribution License (v 1.0)
% This means you may do almost anything with this work of mine, so long as you give me proper credit

Bruk en datamaskin for å navigere til ``Socratic Instrumentation''-nettstedet:

$$\hbox{\tt http://www.ibiblio.org/kuphaldt/socratic/sinst}$$

Når du kommer dit, klikk på lenken for kvartalet (Summer, Fall, Winter eller Spring) du er meldt opp i, og last ned INST200 ``Introduction to Instrumentation'' kursarbeidsark.  Dagens undervisningsøkt vil dekke dag 1 i dette arbeidsarket.

\vskip 10pt

Helt i begynnelsen av dette dokumentet, slik det er for {\it alle} arbeidsarkene for 200-nivå Instrumentation-kurs, finner du en side med tittelen ``How To . . .''.  Finn denne siden og les den grundig, ettersom du vil bli testet på innholdet gjennom hele INST200-kurset.  ``How to . . .''-tipsene refererer til en ``Question 0'' som er en annen side som finnes i hvert kursarbeidsark.  Les også punktene som er listet i Question 0.

%\vskip 10pt

%When the class is finished with the reading, we will review and discuss items in the ``How To . . .'' and ``Question 0'' pages, relating them to the points brainstormed earlier regarding your own professional goals in becoming an instrument technician and the expectations of employers.  This exercise also serves as an introduction showing how all class work is handled in the second year of the Instrumentation program: you complete reading and research assignments prior to each class session, and then you spend nearly the entire class session discussing what you've learned while the instructor challenges you to explore the concepts deeper.  Later in this course you will be quizzed on various details contained in this ``How To . . .'' page.  Feel free to keep a printed copy of this page handy for future INST200 course class sessions, as the quizzes will be open-note!

\vskip 10pt

Instruktøren din vil også dele ut kopier av et samtykkeskjema (``FERPA form'') som du kan signere for å gi tillatelse til å dele dine akademiske resultater med arbeidsgivere.  Dette er frivillig, ikke obligatorisk.  Uten signert samtykke fra studenten forbyr føderal lov enhver instruktør å dele akademiske opplysninger med andre enn studenten og passende ansatte ved høyskolen.

\vskip 10pt

Instruktøren din vil også ha elektroniske kopier (f.eks. minnepinne og/eller CD-ROM) av ``Instrumentation Reference'' tilgjengelig slik at du kan kopiere dem til din personlige datamaskin.  Dette er en samling filer, hovedsakelig hentet fra ulike produsenters nettsteder med deres tillatelse, bestående av veiledninger og rapporter og tekniske manualer som du vil få i oppgave å lese gjennom andreårs-kursene.  Hensikten med denne Referansen er å gi deg rask, frakoblet tilgang slik at du ikke trenger å søke på internett etter disse dokumentene.  Det finnes en fil i rotkatalogen til denne Referansen med navnet ``{\tt 00\_index\_OPEN\_THIS\_FILE.html}'' som du bør åpne i en nettleser.  Hyperlenkene i denne HTML-indeksfilen gjør det mye enklere å finne dokument(ene) du leter etter enn det ville vært å bla gjennom de ulike katalogene i Referansen for å lete etter filnavn.


\vskip 20pt \vbox{\hrule \hbox{\strut \vrule{} {\bf Forslag til sokratisk diskusjon} \vrule} \hrule}

\begin{itemize}
\item{} En av hensiktene med denne øvelsen er å øve på aktive lesestrategier, der du samhandler med teksten for å identifisere og utforske viktige prinsipper.  En effektiv strategi er å skrive ned tanker som dukker opp mens du leser teksten.  Beskriv hvordan denne aktive lesestrategien kan være nyttig i daglige lekseoppgaver.
\item{} For hvert eneste punkt som er listet i ``How To . . .''- og ``Question 0''-sidene, identifiser hvorfor disse punktene er viktige for ditt endelige mål om å bli automatiker/instrumenttekniker.
\item{} Identifiser hvordan kursopplegg og forventninger på INST200-nivå skiller seg fra det du har opplevd tidligere som student, og forklar hvorfor disse forskjellene finnes.
\end{itemize}

\underbar{file i00002}
%(END_QUESTION)





%(BEGIN_ANSWER)


%(END_ANSWER)





%(BEGIN_NOTES)

{\it Du bør skrive ut kopier av ``How To''- og ``Question 0''-sidene (på ett tosidig ark) og ``General Values and Expectations'' på et annet ark for hver eneste student, slik at alle har dem på pulten sin, i tilfelle de ikke har en personlig datamaskin med seg denne dagen.}

\vskip 10pt

De ``Real examples'' viser faktiske spørsmål og scenarier som er blitt stilt til instruktører i Instrumentation-programmet, og kan fungere som utgangspunkt for helklassediskusjoner om hvordan man anvender prinsippene som er listet på ``How To''-siden.  Selv om mange av disse utsagnene faktisk er sammenstillinger av forskjellige utsagn gjort av ulike studenter til ulike tider, er den generelle tematikken og tonen i hver av dem tro mot originalen.

\vskip 10pt

Vær forberedt på å supplere disse eksemplene med live-demonstrasjoner på kateter-PC-en (f.eks. ha karakterregnearket klart til å vise, ha INSTREF klart til å dele osv.).







\vfil \eject
\noindent
{\bf Ekte eksempel:} 

Dagen før eksamen kommer en student bort til instruktøren og spør: ``Hva kommer på morgendagens eksamen?''






\vfil \eject
\noindent
{\bf Ekte eksempel:} 

En student kaller instruktøren bort for hjelp.  Studenten prøver å finne en av instruksjonsmanualene for et bestemt instrument, som er gitt som lesing.  Instruktøren finner studenten mens han blar gjennom de ulike mappene i Instrumentation Reference og leter etter et filnavn som ligner merke og modell på instrumentet som er oppgitt i leksene.  ``Det er så vanskelig å finne noe i denne Instrumentation Reference,'' sier studenten.  ``Finnes det ikke en enklere måte?''






\vfil \eject
\noindent
{\bf Ekte eksempel:} 

Ved starten av et nytt kvartal blir en av de fortsettende andreårsstudentene av instruktøren oppdaget å ha registrert seg til et par feil kurs.  Når instruktøren tar kontakt med studenten for å foreslå at han dropper de feil kursene, sier studenten: ``Personen i registreringen sa at dette var kursene som tilbys dette kvartalet, så jeg meldte meg opp i alle.''






\vfil \eject
\noindent
{\bf Ekte eksempel:} 

En student oppsøker instruktøren midtveis i det akademiske kvartalet.  ``Jeg skulle ønske du la ut karakterer til alle studentene elektronisk, slik de andre instruktørene mine gjør, så vi kan se nøyaktig hvor vi ligger.  Jeg har ingen anelse om hva GPA-en min er akkurat nå, og jeg er bekymret for om jeg kan kvalifisere for et kommende stipend.''






\vfil \eject
\noindent
{\bf Ekte eksempel:} 

En nyutdannet som har vært ute av skolen i 6 måneder sender e-post til instruktøren for å spørre om jobber.  ``Jeg ser etter kontraktarbeid helst i Sør-California, og vet ikke hvor jeg skal begynne å lete.''  Når instruktøren spør hvilke jobber han har søkt på så langt, er svaret: ``Ingen -- jeg er bare i gang.''






\vfil \eject
\noindent
{\bf Ekte eksempel:} 

Under et intervju blir en nyutdannet bedt om å beskrive en situasjon der han måtte løse et sett komplekse problemer for å oppnå et bestemt mål.  Den nyutdannede blir tatt på sengen, vet ikke hva intervjueren ønsker å høre, og føler seg også litt intimidert av spørsmålet fordi han ennå ikke har industriell arbeidserfaring.  Bør han dikte opp en historie som høres bra ut?  Bør han tenke på et hypotetisk scenario der han kan forklare hva han {\it ville} gjort dersom situasjonen oppsto?  Bør han relatere en erfaring fra et prosjekt han jobbet med på skolen, selv om det ikke var en faktisk jobb?






\vfil \eject
\noindent
{\bf Ekte eksempel:} 

En student kommer til timen en dag etter å ha forberedt seg på feil sett med leksespørsmål.  Når han innser dette, forteller han instruktøren om det.  ``Jeg spurte labkameratene mine hvilke spørsmål vi skulle gjøre til i dag, og de ga meg feil informasjon!''







\vfil \eject
\noindent
{\bf Ekte eksempel:} 

En student kommer bort til instruktøren med et spørsmål: ``Foreldrene mine har planlagt ferie den siste uken i oktober, og jeg hadde tenkt å bli med dem.  Går jeg glipp av noe viktig hvis jeg er borte den uken?''





%INDEX% Course organization, career advice
%INDEX% Course organization, course calendar spreadsheet
%INDEX% Course organization, course grading spreadsheet
%INDEX% Course organization, introduce ``How To...'' page
%INDEX% Course organization, Instrumentation Reference
%INDEX% Course organization, job searching
%INDEX% Course organization, second-year course sequence
%INDEX% Course organization, Socratic Instrumentation website
%INDEX% Course organization, student email

%(END_NOTES)
